\subsection{根系的定义}

引理 1. 设 V 是 k 上的向量空间，R 是生成 V 的有限子集。对于每个满足 \(\alpha \in R\)、\(\alpha \neq 0\) 的元素，至多存在一个 V 的反射 s，使得 \(s(\alpha) = -\alpha\) 且 \(s(R) = R\) 。

设 G 为保持 R 不变的 V 的自同构群。由于 R 生成 V，G 同构于 R 的对称群的一个子群，因而是有限的。设 \(s, s'\) 是 V 的反射，满足 \(s(\alpha) = s'(\alpha) = -\alpha\)、\(s(R) = R\)、\(s'(R) = R\)。于是 \(t = ss'\) 属于 G，因而具有有限阶 m。另一方面，

\[
t (\alpha) = \alpha \text { and } t (x) \equiv x \bmod k \alpha \text { for   all } x \in V.
\]

因此，存在一个线性形式 \(f\)，定义在 \(V\) 上，使得

\[
t (x) = x + f (x) \alpha \text {   for   all   } x \in V
\]

且 \(f(\alpha) = 0\)。对 \(n\) 归纳可得

\[
t ^ {n} (x) = x + n f (x) \alpha \text {   for   all   } x \in V.
\]

取 \(n\) 等于 \(m\)，我们看到 \(mf(x) = 0\) 对所有 \(x \in V\) 成立，所以 \(f = 0\)、\(t = 1\) 且 \(s = s'\)。

定义 1. 设 V 是 k 上的向量空间，R 是 V 的子集。若满足下列条件，则称 R 是 V 中的一个根系：

(RS \(_{I}\) ) R 是有限的，不含 0，且生成 V。

(RS \(_{II}\) ) 对所有 \(\alpha \in \mathbb{R}\)，存在元素 \(\alpha^{\prime}\)，属于 V 的对偶空间 V \(^{*}\)，使得 \(\langle \alpha, \alpha^{\prime} \rangle = 2\)，并且反射 \(s_{\alpha, \alpha'}\)（cf.~Chap. V, § 2）保持 R 不变。

(RS \(_{III}\) ) 对所有 \(\alpha \in R\)，有 \(\alpha^{\sim}(R) \subseteq \mathbf{Z}\)。

由引理 1，反射 \(s_{\alpha, \alpha'}\)（从而线性形式 \(\alpha'\) 也）由 \(\alpha\) 唯一确定，因此 \((\mathrm{RS}_{\mathrm{III}})\) 有意义。置 \(s_{\alpha, \alpha'} = s_{\alpha}\)。于是 \(s_{\alpha}(x) = x - \langle \alpha', x \rangle \alpha\)，对所有 \(x \in V\) 成立。

R 的元素称为（所考虑系统的）根。V 的维数称为该系统的秩。

保持 R 不变的 V 的自同构称为 R 的自同构。它们构成有限群，记为 A(R)。由 \(s_{\alpha}\) 生成的 A(R) 子群称为 R 的 Weyl 群，记为 W(R)，或简记为 W。

注。1）设 \(k'\) 是 \(k\) 的扩域。将 V 典范地视为 \(V \otimes k'\) 的子集，将 \(V^*\) 视为 \(V^* \otimes k' = (V \otimes k')^*\) 的子集。则 R 是 \(V \otimes k'\) 中的根系，且 \(\alpha'\) 与此前相同。

引理 2. 设 R 是 V 中的根系。设 \((x|y)\) 是 V 上的非退化、在 W(R) 下不变的对称双线性型。利用此型将 V 与 \(V^{*}\) 等同。若 \(\alpha \in R\)，则 \(\alpha\) 非等方，且

\[
\alpha^ {\prime} = \frac {2 \alpha}{(\alpha | \alpha)}.
\]

这由 Chap. V, § 2, no. 3 的公式 (4) 得出。

命题 1. 设 \(\mathrm{V}_{\mathbf{Q}}\)（相应地，\(\mathrm{V}_{\mathbf{Q}}^{*}\)）是一个 \(\mathbf{Q}\) -向量子空间，由 \(\mathrm{V}\)（相应地，\(\mathrm{V}^{*}\)）中的 \(\alpha\)（相应地，\(\alpha^{-}\)）生成。于是 \(\mathrm{V}_{\mathbf{Q}}\)（相应地，\(\mathrm{V}_{\mathbf{Q}}^{*}\)）是 \(\mathbf{Q}\) -结构，定义在 \(\mathrm{V}\)（相应地，\(\mathrm{V}^{*}\)）上（Algebra, Chap. II, § 8, no. 1）。将典范双线性型限制在 \(\mathrm{V}_{\mathbf{Q}} \times \mathrm{V}_{\mathbf{Q}}^{*}\) 上（该型定义在 \(\mathrm{V} \times \mathrm{V}^{*}\) 上），得到把空间 \(\mathrm{V}_{\mathbf{Q}}, \mathrm{V}_{\mathbf{Q}}^{*}\) 各自与另一个空间的对偶等同的映射。集合 \(\mathbb{R}\) 是 \(\mathrm{V}_{\mathbf{Q}}\) 中的根系。

若 \(k = \mathbf{R}\) ，则存在 \(V\) 上在 \(W(R)\) 下不变的内积（Integration, Chap. VII, § 3, no. 1, Prop. 1）；引理 2 现在表明 \(\alpha^{-}\) 生成 \(V^{*}\) 。由注记 1，\(\alpha^{-}\) 同样生成 \(V^{*}\)（当 \(k = \mathbf{Q}\) 时）。现在转到一般情形。置 \(E = V_{\mathbf{Q}}\) 。由 \((RS_{III})\) ，每个 \(\alpha^{-}\) 将 \(E\) 映到 \(\mathbf{Q}\) ，因而定义元素 \(\tilde{\alpha}\) 为 \(E^{*}\) 中的一个元素。显然 \(R\) 是 \(E\) 中的根系，并且与 \(\alpha\) 对应的 \(E^{*}\) 中的元素是 \(\tilde{\alpha}\) 。由上文所述，\(\tilde{\alpha}\) 生成向量空间 \(E^{*}\) 。考虑典范同态 \(i: E \otimes_{\mathbf{Q}} k \to V\) 及其转置 \({}^{t}i: V^{*} \to E^{*} \otimes_{\mathbf{Q}} k\) 。由于 \(R\) 生成 \(V\) ，\({}^{t}i\) 是单射；但 \({}^{t}i\) 的像包含 \(\tilde{\alpha}\) ，所以 \({}^{t}i\) 是满射。由此最终得到 \(i\) 和 \({}^{t}i\) 都是同构。因而可以把 \(V\) 与 \(E \otimes k\) 等同，把 \(V^{*}\) 与 \(E^{*} \otimes k\) 等同，把 \(\alpha^{-}\) 与 \(\tilde{\alpha}\) 等同，并把 \(V_{\mathbf{Q}}^{*}\) 与 \(E^{*}\) 等同。于是，\(V_{\mathbf{Q}}\)（相应地，\(V_{\mathbf{Q}}^{*}\)）是 \(Q\) -结构，定义在 \(V\)（相应地，\(V^{*}\)）上。将典范双线性型限制在 \(V_{\mathbf{Q}} \times V_{\mathbf{Q}}^{*}\) 上（该型定义在 \(V \times V^{*}\) 上）后，可以与 \(E \times E^{*}\) 上的典范双线性型等同，命题得证。

注记。2）命题 1 将根系的研究归结为 \(k = \mathbf{Q}\) 的情形。注 1 又将其进一步归结为实向量空间 \(V_{\mathbf{R}} = V_{\mathbf{Q}} \otimes_{\mathbf{Q}} \mathbf{R}\) 中的根系研究。与这些不同系统相关的 Weyl 群可典范地等同。

\begin{enumerate}
\def\labelenumi{\arabic{enumi})}
\setcounter{enumi}{2}
\tightlist
\item
  由于 \(\alpha^{\sim}\) 生成 \(V^{*}\)，群 \(W(R)\)（视为 \(\mathbf{GL}(V_{\mathbf{R}})\) 的子群）是本质的（Chap. V, § 3, no. 7）。此外，Chap. V, § 3, no. 2 的 Th. 1 的 Cor. 表明，属于 \(W(R)\) 的唯一反射是 \(s_{\alpha}\)。
\end{enumerate}

命题 2. \(\alpha^{\vee}\) 在 \(V^{*}\) 中构成根系，并且 \(\alpha^{\vee} = \alpha\) 对所有 \(\alpha \in R\) 成立。

由命题 1，\(\alpha^{\prime}\) 满足 \((\mathrm{RS}_{\mathrm{I}})\)。由于 \(s_{\alpha, \alpha'}\) 是赋予子集 R 的向量空间 V 的自同构，\({}^{t}(s_{\alpha, \alpha'})^{-1}\) 保持由 \(\alpha'\) 构成的集合 R' 不变；但 \({}^{t}(s_{\alpha, \alpha'})^{-1} = s_{\alpha', \alpha}\)，这证明了 R' 满足 \((\mathrm{RS}_{\mathrm{II}})\)，且 \(\alpha' = \alpha\)。最后，\(\langle \alpha', \beta \rangle \in \mathbf{Z}\) 对所有 \(\alpha' \in \mathrm{R'}\) 和 \(\beta \in \mathrm{R}\) 成立，所以 R' 满足 \((\mathrm{RS}_{\mathrm{III}})\)。

集合 \(R^{-}\) 称为 R 的逆根系。映射 \(\alpha \mapsto \alpha^{-}\) 是从 R 到 \(R^{-}\) 的双射，称为从 R 到 \(R^{-}\) 的典范双射。注意，若 \(\alpha, \beta\) 是 R 的元素且 \(\alpha + \beta \in R\)，则一般而言 \((\alpha + \beta)^{-} \neq \alpha^{-} + \beta^{-}\)。

由于 \(s_{\alpha}(\alpha) = -\alpha\)，公理 (RSII) 表明 \(-R = R\)。显然 \((- \alpha)^{\sim} = -\alpha^{\sim}\)，且 \(-1 \in A(R)\)（但 \(-1 \in W(R)\) 并不总成立）。

等式 \({}^{t}(s_{\alpha,\alpha^{-}})^{-1}=s_{\alpha^{-},\alpha}\) 表明映射 \(u\mapsto{}^{t}u^{-1}\) 是从群 W(R) 到群 W(R \(^{*}\) ) 的同构。我们借助此同构将这两个群等同；换言之，认为 W(R) 同时作用于 V 和 V \(^{*}\)。A(R) 同理。

命题 3. 对 \(x, y \in V\) ，置

\[
(x | y) = \sum_ {\alpha \in R} \langle \alpha^ {\vee}, x \rangle \langle \alpha^ {\vee}, y \rangle .
\]

则 \((x|y)\) 是 \(\mathbf{V}\) 上的非退化对称双线性型，在 \(\mathbf{A}(\mathbb{R})\) 下不变。对 \(x, y \in \mathrm{V}_{\mathbf{Q}}\)，有 \((x|y) \in \mathbf{Q}\)。\((x|y)\) 的典范延拓到

\[
V _ {R} = V _ {Q} \otimes_ {Q} R
\]

是非退化且正定的。

显然 \((x|y)\) 是 V 上的对称双线性型。若 \(g\in \mathrm{A}(\mathbb{R})\)

\[
(g (x) | g (y)) = \sum_ {\alpha \in R} \left\langle {} ^ {t} g \left(\alpha^ {-}\right), x \right\rangle \left\langle {} ^ {t} g \left(\alpha^ {-}\right), y \right\rangle = (x | y)
\]

因为 \((^t g)(\mathbb{R}^\sim) = \mathbb{R}^\sim\)。若 \(x, y \in \mathrm{V}_{\mathbf{Q}}\)，则 \((x|y) \in \mathbf{Q}\) 由 \((\mathrm{RS}_{\mathrm{III}})\) 得出。若 \(z \in \mathrm{V}_{\mathbf{R}}\)，则 \((z|z) = \sum_{\alpha \in \mathbb{R}} \langle \alpha^\sim, z \rangle^2 \geqslant 0\)，且有 \((z|z) > 0\)，只要 \(z \neq 0\)（由命题 1），所以 \((x|y)\) 到 \(V_{R}\) 的典范延拓是正定且非退化的。因此，\((x|y)\) 限制到 \(V_{Q}\) 后是非退化的，从而 V 上的型 \((x|y)\) 是非退化的。

命题 4. (i) 设 X 是 R 的子集，设 \(V_{X}\) 是由 X 生成的 V 的向量子空间，设 \(V'_{X}\) 是 \(V^{*}\) 中的向量子空间，由 \(\alpha'\) 生成，其中 \(\alpha \in X\)。于是 V 是 \(V_{X}\) 与 \(V'_{X}\) 的正交补的直和，\(V^{*}\) 是 \(V'_{X}\) 与 \(V_{X}\) 的正交补的直和，并且 \(V'_{X}\) 与 \(V_{X}\) 的对偶等同。

\begin{enumerate}
\def\labelenumi{(\roman{enumi})}
\setcounter{enumi}{1}
\tightlist
\item
  \(R \cap V_{X}\) 是 \(V_{X}\) 中的根系，并且从 \(R \cap V_{X}\) 到其逆根系的典范双射等同于将映射 \(\alpha \mapsto \alpha^{-}\) 限制到 \(R \cap V_{X}\) 所得的映射。
\end{enumerate}

由注记 2，可假设 \(k = \mathbf{R}\)。利用命题 3 的对称双线性型将 \(V\) 与 \(V^*\) 等同。由引理 2，有 \(\alpha^{-} = \frac{2\alpha}{(\alpha|\alpha)}\) 对所有 \(\alpha \in R\) 成立（引理 2）。\(V\) 的每个向量子空间都是非等方的，命题现已显然。

推论。设 \(V_{1}\) 是 V 的一个向量子空间，设 \(V_{2}\) 是由 \(R \cap V_{1}\) 生成的向量子空间。则 \(R \cap V_{1}\) 是 \(V_{2}\) 中的根系。

这由将 (ii) 应用于 \(\mathrm{X} = \mathrm{R} \cap \mathrm{V}_1\) 得出。

对 \(\alpha, \beta \in \mathbb{R}\)，置

\[
\langle \alpha , \beta^ {\prime} \rangle = n (\alpha , \beta).\tag{1}
\]

于是

\[
n (\alpha , \alpha) = 2\tag{2}
\]

\[
n (- \alpha , \beta) = n (\alpha , - \beta) = - n (\alpha , \beta).\tag{3}
\]

由 \((\mathrm{RS}_{\mathrm{III}})\)

\[
n (\alpha , \beta) \in \mathbf {Z}.\tag{4}
\]

由 \(n(\alpha, \beta)\) 的定义，

\[
s _ {\beta} (\alpha) = \alpha - n (\alpha , \beta) \beta .\tag{5}
\]

公式 (1) 和命题 2 蕴含

\[
n (\alpha , \beta) = n \left(\beta^ {-}, \alpha^ {-}\right).\tag{6}
\]

设 \((x|y)\) 是 V 上的非退化、在 W(R) 下不变的对称双线性型（命题 3）。由引理 2，

\[
n (\alpha , \beta) = \frac {2 (\alpha | \beta)}{(\beta | \beta)}.\tag{7}
\]

于是

\[
n (\alpha , \beta) = 0 \Longleftrightarrow n (\beta , \alpha) = 0 \Longleftrightarrow (\alpha | \beta) = 0 \Longleftrightarrow s _ {\alpha} \text {   and   } s _ {\beta} \text {   commute.   (8) }
\]

\[
\text {   If   } (\alpha | \beta) \neq 0, \text {   then   } \frac {n (\beta , \alpha)}{n (\alpha , \beta)} = \frac {(\beta | \beta)}{(\alpha | \alpha)}.\tag{9}
\]

